\documentclass[letterpaper]{article}
\usepackage[submission]{aaai2027}
\usepackage[hyphens]{url}
\usepackage{graphicx}
\urlstyle{rm}
\def\UrlFont{\rm}
\usepackage{natbib}
\usepackage{caption}
\usepackage{booktabs}
\usepackage{amsmath}
\usepackage{amssymb}
\frenchspacing
\pdfinfo{/TemplateVersion (2027.1)}
\setcounter{secnumdepth}{2}
\renewcommand{\thesection}{S\arabic{section}}
\renewcommand{\thetable}{S\arabic{table}}
\title{Supplementary Material:\\Sensor Metadata and Axis-Aware Tokenization Shape Zero-Shot Reconstruction in OPM-MEG}
\author{Anonymous Authors}
\affiliations{}
% Generated by factorial_results.py. Do not transcribe values.
\newcommand{\FxRawMagMonolithicMean}{0.3726}
\newcommand{\FxRawMagMonolithicCI}{[0.3409, 0.4047]}
\newcommand{\FxRawMagMonolithicCount}{10/10}
\newcommand{\FxRawMagMonolithicP}{0.002}
\newcommand{\FxRawMagRandomMean}{0.4019}
\newcommand{\FxRawMagRandomCI}{[0.3721, 0.4316]}
\newcommand{\FxRawMagRandomCount}{10/10}
\newcommand{\FxRawMagRandomP}{0.002}
\newcommand{\FxRawMagAxisPartitionedMean}{0.5138}
\newcommand{\FxRawMagAxisPartitionedCI}{[0.4846, 0.5441]}
\newcommand{\FxRawMagAxisPartitionedCount}{10/10}
\newcommand{\FxRawMagAxisPartitionedP}{0.002}
\newcommand{\FxRawGradMonolithicMean}{0.3728}
\newcommand{\FxRawGradMonolithicCI}{[0.3388, 0.4064]}
\newcommand{\FxRawGradMonolithicCount}{10/10}
\newcommand{\FxRawGradMonolithicP}{0.002}
\newcommand{\FxRawGradRandomMean}{0.4174}
\newcommand{\FxRawGradRandomCI}{[0.3890, 0.4454]}
\newcommand{\FxRawGradRandomCount}{10/10}
\newcommand{\FxRawGradRandomP}{0.002}
\newcommand{\FxRawGradAxisPartitionedMean}{0.4862}
\newcommand{\FxRawGradAxisPartitionedCI}{[0.4652, 0.5066]}
\newcommand{\FxRawGradAxisPartitionedCount}{10/10}
\newcommand{\FxRawGradAxisPartitionedP}{0.002}
\newcommand{\FxSofthfcMagMonolithicMean}{0.3456}
\newcommand{\FxSofthfcMagMonolithicCI}{[0.3102, 0.3807]}
\newcommand{\FxSofthfcMagMonolithicCount}{10/10}
\newcommand{\FxSofthfcMagMonolithicP}{0.002}
\newcommand{\FxSofthfcMagRandomMean}{0.3785}
\newcommand{\FxSofthfcMagRandomCI}{[0.3457, 0.4105]}
\newcommand{\FxSofthfcMagRandomCount}{10/10}
\newcommand{\FxSofthfcMagRandomP}{0.002}
\newcommand{\FxSofthfcMagAxisPartitionedMean}{0.4856}
\newcommand{\FxSofthfcMagAxisPartitionedCI}{[0.4512, 0.5201]}
\newcommand{\FxSofthfcMagAxisPartitionedCount}{10/10}
\newcommand{\FxSofthfcMagAxisPartitionedP}{0.002}
\newcommand{\FxSofthfcGradMonolithicMean}{0.3471}
\newcommand{\FxSofthfcGradMonolithicCI}{[0.3141, 0.3796]}
\newcommand{\FxSofthfcGradMonolithicCount}{10/10}
\newcommand{\FxSofthfcGradMonolithicP}{0.002}
\newcommand{\FxSofthfcGradRandomMean}{0.3978}
\newcommand{\FxSofthfcGradRandomCI}{[0.3702, 0.4247]}
\newcommand{\FxSofthfcGradRandomCount}{10/10}
\newcommand{\FxSofthfcGradRandomP}{0.002}
\newcommand{\FxSofthfcGradAxisPartitionedMean}{0.4723}
\newcommand{\FxSofthfcGradAxisPartitionedCI}{[0.4500, 0.4944]}
\newcommand{\FxSofthfcGradAxisPartitionedCount}{10/10}
\newcommand{\FxSofthfcGradAxisPartitionedP}{0.002}
\newcommand{\FxRawMagGroupingMean}{+0.1120}
\newcommand{\FxRawMagGroupingCI}{[+0.1048, +0.1208]}
\newcommand{\FxRawMagGroupingCount}{10/10}
\newcommand{\FxRawMagGroupingP}{0.002}
\newcommand{\FxRawMagMonoToRandomMean}{+0.0292}
\newcommand{\FxRawMagMonoToRandomCI}{[+0.0262, +0.0322]}
\newcommand{\FxRawMagMonoToRandomCount}{10/10}
\newcommand{\FxRawMagMonoToRandomP}{0.002}
\newcommand{\FxRawGradGroupingMean}{+0.0688}
\newcommand{\FxRawGradGroupingCI}{[+0.0567, +0.0805]}
\newcommand{\FxRawGradGroupingCount}{10/10}
\newcommand{\FxRawGradGroupingP}{0.002}
\newcommand{\FxRawGradMonoToRandomMean}{+0.0446}
\newcommand{\FxRawGradMonoToRandomCI}{[+0.0390, +0.0504]}
\newcommand{\FxRawGradMonoToRandomCount}{10/10}
\newcommand{\FxRawGradMonoToRandomP}{0.002}
\newcommand{\FxSofthfcMagGroupingMean}{+0.1071}
\newcommand{\FxSofthfcMagGroupingCI}{[+0.1002, +0.1143]}
\newcommand{\FxSofthfcMagGroupingCount}{10/10}
\newcommand{\FxSofthfcMagGroupingP}{0.002}
\newcommand{\FxSofthfcMagMonoToRandomMean}{+0.0328}
\newcommand{\FxSofthfcMagMonoToRandomCI}{[+0.0292, +0.0363]}
\newcommand{\FxSofthfcMagMonoToRandomCount}{10/10}
\newcommand{\FxSofthfcMagMonoToRandomP}{0.002}
\newcommand{\FxSofthfcGradGroupingMean}{+0.0745}
\newcommand{\FxSofthfcGradGroupingCI}{[+0.0656, +0.0833]}
\newcommand{\FxSofthfcGradGroupingCount}{10/10}
\newcommand{\FxSofthfcGradGroupingP}{0.002}
\newcommand{\FxSofthfcGradMonoToRandomMean}{+0.0507}
\newcommand{\FxSofthfcGradMonoToRandomCI}{[+0.0448, +0.0564]}
\newcommand{\FxSofthfcGradMonoToRandomCount}{10/10}
\newcommand{\FxSofthfcGradMonoToRandomP}{0.002}
\newcommand{\FxRawTagMonolithicMean}{+0.0002}
\newcommand{\FxRawTagMonolithicCI}{[-0.0104, +0.0106]}
\newcommand{\FxRawTagMonolithicCount}{5/10}
\newcommand{\FxRawTagMonolithicP}{0.969}
\newcommand{\FxRawTagRandomMean}{+0.0155}
\newcommand{\FxRawTagRandomCI}{[+0.0065, +0.0234]}
\newcommand{\FxRawTagRandomCount}{9/10}
\newcommand{\FxRawTagRandomP}{0.014}
\newcommand{\FxRawTagAxisPartitionedMean}{-0.0277}
\newcommand{\FxRawTagAxisPartitionedCI}{[-0.0408, -0.0150]}
\newcommand{\FxRawTagAxisPartitionedCount}{1/10}
\newcommand{\FxRawTagAxisPartitionedP}{0.006}
\newcommand{\FxRawInteractionMean}{-0.0432}
\newcommand{\FxRawInteractionCI}{[-0.0548, -0.0329]}
\newcommand{\FxRawInteractionCount}{0/10}
\newcommand{\FxRawInteractionP}{0.002}
\newcommand{\FxSofthfcTagMonolithicMean}{+0.0014}
\newcommand{\FxSofthfcTagMonolithicCI}{[-0.0108, +0.0120]}
\newcommand{\FxSofthfcTagMonolithicCount}{6/10}
\newcommand{\FxSofthfcTagMonolithicP}{0.824}
\newcommand{\FxSofthfcTagRandomMean}{+0.0193}
\newcommand{\FxSofthfcTagRandomCI}{[+0.0083, +0.0287]}
\newcommand{\FxSofthfcTagRandomCount}{9/10}
\newcommand{\FxSofthfcTagRandomP}{0.012}
\newcommand{\FxSofthfcTagAxisPartitionedMean}{-0.0133}
\newcommand{\FxSofthfcTagAxisPartitionedCI}{[-0.0276, +0.0008]}
\newcommand{\FxSofthfcTagAxisPartitionedCount}{5/10}
\newcommand{\FxSofthfcTagAxisPartitionedP}{0.115}
\newcommand{\FxSofthfcInteractionMean}{-0.0326}
\newcommand{\FxSofthfcInteractionCI}{[-0.0431, -0.0233]}
\newcommand{\FxSofthfcInteractionCount}{0/10}
\newcommand{\FxSofthfcInteractionP}{0.002}
\newcommand{\FxRawMagMinimum}{0.1009}
\newcommand{\FxRawMagMaximum}{0.1215}
\newcommand{\FxRawMagSd}{0.0049}
\newcommand{\FxRawGradMinimum}{0.0591}
\newcommand{\FxRawGradMaximum}{0.0781}
\newcommand{\FxRawGradSd}{0.0049}
\newcommand{\FxSofthfcMagMinimum}{0.0969}
\newcommand{\FxSofthfcMagMaximum}{0.1178}
\newcommand{\FxSofthfcMagSd}{0.0049}
\newcommand{\FxSofthfcGradMinimum}{0.0638}
\newcommand{\FxSofthfcGradMaximum}{0.0836}
\newcommand{\FxSofthfcGradSd}{0.0051}
\newcommand{\FxWeightingBound}{0.0009}
\newcommand{\FxRawMagMonoToRandomPercent}{20.7}
\newcommand{\FxRawMagGroupingPercent}{79.3}
\newcommand{\FxRawGradMonoToRandomPercent}{39.3}
\newcommand{\FxRawGradGroupingPercent}{60.7}
\newcommand{\FxSofthfcMagMonoToRandomPercent}{23.5}
\newcommand{\FxSofthfcMagGroupingPercent}{76.5}
\newcommand{\FxSofthfcGradMonoToRandomPercent}{40.5}
\newcommand{\FxSofthfcGradGroupingPercent}{59.5}
\newcommand{\FxProbeMinimum}{47.53}
\newcommand{\FxProbeMaximum}{51.88}
\newcommand{\FxClassicalRawMean}{53.54}
\newcommand{\FxClassicalRawCI}{[49.93, 56.75]}
\newcommand{\FxClassicalSofthfcMean}{54.50}
\newcommand{\FxClassicalSofthfcCI}{[51.03, 57.65]}
\newcommand{\SaZeroOrientationCount}{21/21}
\newcommand{\SaFlipOrientationCount}{21/21}
\newcommand{\SaScrambleChannelsCount}{19/21}

\begin{document}
\maketitle

\noindent
This supplement presents additional methods, participant-level reconstruction results,
paired comparisons, decoding results, and dataset details.

\section{Soft Homogeneous Field Correction}

\subsection{Setup}
Following \citet{tierney2021hfc}, we approximate distant magnetic interference as a spatially uniform field $\mathbf{B}_{\mathrm{ext}} \in \mathbb{R}^{3}$. A sensor $c$ with unit
sensitive orientation $\mathbf{u}_{c}$ measures the projection
$m_{c} = \mathbf{u}_{c}^{\top}\mathbf{B}_{\mathrm{ext}}$, so across an array of $C$
channels the induced measurement is $\mathbf{m}_{\mathrm{ext}} = \mathbf{Q}\mathbf{B}_{\mathrm{ext}}$,
where row $c$ of $\mathbf{Q} \in \mathbb{R}^{C \times 3}$ is $\mathbf{u}_{c}^{\top}$.

Conventional homogeneous field correction \citep{tierney2021hfc} projects onto the null space of
$\mathbf{Q}^{\top}$:
\[
\mathbf{P}_{\mathrm{hard}} = \mathbf{I}_{C} - \mathbf{Q}(\mathbf{Q}^{\top}\mathbf{Q})^{-1}\mathbf{Q}^{\top}.
\]
When $\mathbf{Q}$ has full column rank, $\mathrm{rank}(\mathbf{P}_{\mathrm{hard}}) = C - 3$, making the projected channel covariance rank-deficient.

\subsection{Soft variant}
The cited method uses full projection. Here we evaluate a soft shrinkage variant parameterized by $\alpha \in [0,1)$:
\[
\mathbf{P}_{\mathrm{soft}}(\alpha) = \mathbf{I}_{C} - \alpha\,\mathbf{Q}_{R}\mathbf{Q}_{R}^{\top},
\]
where $\mathbf{Q} = \mathbf{Q}_{R}\mathbf{R}$ is a thin QR decomposition and
$\mathbf{Q}_{R}^{\top}\mathbf{Q}_{R} = \mathbf{I}_{3}$. Its eigenvalues are $1-\alpha$
on the three directions spanning $\mathrm{col}(\mathbf{Q})$ and $1$ on the remaining
$C-3$, so for any $\alpha < 1$ the operator has full rank $C$.

The coefficient $\alpha = 0.25$ is an analysis choice here, not a value established by \citet{tierney2021hfc}. At this value, components lying in the uniform-field subspace
are scaled by $0.75$ in amplitude, a power reduction of $1 - 0.75^{2} = 43.75\%$.

\subsection{What this does to neural signal}
A neural component orthogonal to $\mathrm{col}(\mathbf{Q})$ satisfies
$\mathbf{Q}_{R}^{\top}\mathbf{s} = \mathbf{0}$ and is returned unchanged. This is not a
claim that neural signal is preserved. A dipolar field pattern is not in general
orthogonal to the uniform-field subspace, and whatever part of it does project onto
$\mathrm{col}(\mathbf{Q})$ receives the same $25\%$ amplitude shrinkage as the
interference. Soft-HFC is therefore a lossy operation on signal as well as noise,
which is why the main paper holds filtering constant within every contrast rather
than treating it as a neutral preprocessing step.

\section{Single-Axis Metadata Interventions}

All conditions share channel selection, trials and Soft-HFC. Only the metadata handed
to the tokenizer differs, so any difference between columns of
Table~\ref{tab:s-single-per-subject} is a property of the learned embedding rather than
of the recording. Table~\ref{tab:s-single-contrasts} gives the paired tests.

The dataset documentation states that FACE-DEC sensor coordinates were approximated
from an Elekta layout for topographic display and are not measured OPM positions. The
geometry rows should therefore be read as sensitivity to supplied metadata, not as
evidence about forward-field accuracy.

\section{Tag by Partition Factorial}
The experiment crosses MAG/GRAD tags, raw/Soft-HFC inputs and monolithic/random/axis partitions. Random group sizes match each participant's $Z$, $X$ and $Y$ counts exactly in the generating script, with seeds 42 and 101--109. Saved group assignments are unavailable for an independent metadata-level check. We average seeds within participant before paired inference, retaining ten independent participants. Tables~\ref{tab:s-factorial-cohort}--\ref{tab:s-factorial-seeds} report the full-array reconstruction results.

PCC is averaged over channels and temporal windows within each trial, then over trials within participant. Each participant contributes equally to cohort estimates.

The tag effect is configuration-dependent. The raw monolithic contrast is $\FxRawTagMonolithicMean$ ($\FxRawTagMonolithicCI$, $p=\FxRawTagMonolithicP$), while axis grouping retains its advantage under both tags. The paired change in the grouping advantage under GRAD is $\FxRawInteractionMean$ raw and $\FxSofthfcInteractionMean$ with Soft-HFC, both negative in every participant. This interaction does not by itself identify a shared mechanism. Under Soft-HFC, the monolithic-to-random contrast accounts for $\FxSofthfcMagMonoToRandomPercent\%$ under MAG and $\FxSofthfcGradMonoToRandomPercent\%$ under GRAD of the respective axis-over-monolithic gains. These are descriptive partition components: random partitioning changes attention composition as well as query budget, so a capacity contribution cannot be isolated. Tests are exploratory and unadjusted for multiple comparisons.

\section{Decoding}
The classical decoder uses scikit-learn \texttt{LogisticRegression} with $L_2$ penalty at the library default $C=1.0$, fixed for every configuration and never tuned on a held-out fold, so no selection leakage is possible. The pretrained probe uses the same solver and folds. The sensor-type codes are checkpoint-specific integers, \texttt{MAG=1} and \texttt{GRAD=2} in the released tiny configuration, and any reuse of this intervention must read them from the checkpoint rather than assume them.


Table~\ref{tab:s-factorial-decoding} gives all twelve full-array configurations and three decoders, evaluated with five-fold cross-validation within participant. Random-partition decoding uses seed 42 with group sizes matched to the axes, while reconstruction averages ten seeds. Every pretrained-probe interval includes chance. Table~\ref{tab:s-decoding-contrasts} gives the paired axis-minus-monolithic and axis-minus-random probe contrasts; every interval includes zero. These are exploratory tests without multiplicity correction, not evidence of equivalence.

The classical features are six log band powers spanning 1--96\,Hz plus four temporal window means per channel. They depend on filtering but not on sensor tag or tokenizer partition. The six classical rows within each filtering regime therefore share one underlying measurement. The Soft-HFC interval excludes chance, while the unfiltered interval does not. The classical cohort mean exceeds the pretrained probe in every configuration, but this ordering alone does not establish a significant paired advantage or rule out a task-level noise floor.

\section{Model Architecture}

The pretrained BrainOmni tiny configuration is evaluated with frozen weights.
The tokenizer conditions on sensor position, unit orientation, and a discrete
sensor type, and uses residual vector quantization \citep{xiao2025brainomni}. In the evaluated configuration, the sequence backbone is
bidirectional (\texttt{causal=False}) and uses masked-token prediction over
codebook entries with a mask ratio of $0.25$. Attention operates within trials,
not across trials. The number of learned spatial queries is loaded from the
checkpoint.

\noindent
Axis partitioning is an inference-time reorganization only. The $X$, $Y$ and $Z$ groups
pass separately through the same frozen tokenizer and the reconstructions are
reassembled, giving three spatial-attention passes instead of one. No weights change.

\section{Dataset Details}
FACE-DEC, OpenNeuro \texttt{ds005107} v1.0.0, DOI \texttt{10.18112/openneuro.ds005107.v1.0.0},
21 participants and the 50 channels matched across participants by rounded position and
orientation metadata. Triaxial OPM, Zenodo record \texttt{10.5281/zenodo.11126593}, the
10 adult participants, 84 cached trials each, 155 to 176 retained channels. SQUID control,
OpenNeuro \texttt{ds000117} v1.0.4, participant \texttt{sub-01}, run \texttt{01}.

\begin{table*}[htbp]\centering\footnotesize
\setlength{\tabcolsep}{2.5pt}
\begin{tabular}{lrrrrr}
\toprule
Participant & Default & GRAD tag & Zero ori. & Rev. ori. & Scrambled \\
\midrule
P01 & 0.2717 & 0.3950 & 0.2600 & 0.2563 & 0.2327 \\
P02 & 0.2522 & 0.3735 & 0.2413 & 0.2418 & 0.2331 \\
P03 & 0.2435 & 0.3713 & 0.2299 & 0.2268 & 0.2265 \\
P04 & 0.2444 & 0.3976 & 0.2310 & 0.2250 & 0.2420 \\
P05 & 0.2484 & 0.3944 & 0.2386 & 0.2445 & 0.2466 \\
P06 & 0.2573 & 0.3910 & 0.2461 & 0.2439 & 0.2279 \\
P07 & 0.2580 & 0.3940 & 0.2501 & 0.2511 & 0.2350 \\
P08 & 0.2630 & 0.3890 & 0.2518 & 0.2502 & 0.2327 \\
P09 & 0.2641 & 0.3736 & 0.2549 & 0.2504 & 0.2382 \\
P10 & 0.2486 & 0.3708 & 0.2352 & 0.2306 & 0.2376 \\
P11 & 0.2546 & 0.3855 & 0.2384 & 0.2384 & 0.2300 \\
P12 & 0.2415 & 0.3839 & 0.2270 & 0.2241 & 0.2355 \\
P13 & 0.2302 & 0.3680 & 0.2134 & 0.2063 & 0.2347 \\
P14 & 0.2579 & 0.3864 & 0.2431 & 0.2425 & 0.2352 \\
P15 & 0.2352 & 0.3753 & 0.2159 & 0.2109 & 0.2251 \\
P16 & 0.2571 & 0.3836 & 0.2462 & 0.2446 & 0.2350 \\
P17 & 0.2564 & 0.3884 & 0.2467 & 0.2425 & 0.2251 \\
P18 & 0.2509 & 0.3745 & 0.2327 & 0.2307 & 0.2303 \\
P19 & 0.2562 & 0.3813 & 0.2363 & 0.2370 & 0.2268 \\
P20 & 0.2373 & 0.3771 & 0.2163 & 0.2137 & 0.2258 \\
P21 & 0.2226 & 0.3647 & 0.1990 & 0.2020 & 0.2368 \\
\midrule
Mean & 0.2501 & 0.3819 & 0.2359 & 0.2340 & 0.2330 \\
\bottomrule\end{tabular}
\caption{Zero-shot reconstruction PCC per participant on FACE-DEC (ds005107), 50 shared channels, Soft-HFC $\alpha=0.25$. Participants are indexed in the saved record order, which is not the dataset's own participant numbering. Waveforms are identical across columns; only the metadata supplied to the tokenizer differs.}
\label{tab:s-single-per-subject}\end{table*}

\begin{table*}[htbp]\centering\small
\begin{tabular}{lrrrr}
\toprule
Contrast vs default & Mean $\Delta$ & 95\% CI & Consistent & Exact $p$ \\
\midrule
Gradiometer tag & $+0.1318$ & $[+0.1275, +0.1360]$ & 21/21 & 9.54e-07 \\
Zero orientation & $-0.0142$ & $[-0.0160, -0.0124]$ & 21/21 & 9.54e-07 \\
Reversed orientation & $-0.0161$ & $[-0.0183, -0.0139]$ & 21/21 & 9.54e-07 \\
Scrambled geometry & $-0.0171$ & $[-0.0225, -0.0115]$ & 19/21 & 4.58e-05 \\
\bottomrule\end{tabular}
\caption{Paired single-axis contrasts, 21 participants. Confidence intervals are percentile bootstrap over participants, 20{,}000 resamples, seed 0. The consistency column counts participants moving in the stated direction. $p$ is the exact two-sided sign-flip permutation value enumerating all $2^{21}$ assignments.}
\label{tab:s-single-contrasts}\end{table*}

\begin{table*}[htbp]\centering\small
\begin{tabular}{llrr}
\toprule
Configuration & HFC & Mean PCC & 95\% CI \\
\midrule
Monolithic & no & 0.3726 & $[0.3405, 0.4048]$ \\
Monolithic & Soft & 0.3461 & $[0.3100, 0.3815]$ \\
Random 3-way & no & 0.3995 & $[0.3689, 0.4286]$ \\
Random 3-way & Soft & 0.3753 & $[0.3414, 0.4074]$ \\
Axis partition & no & 0.5130 & $[0.4832, 0.5436]$ \\
Axis partition & Soft & 0.4849 & $[0.4500, 0.5197]$ \\
\midrule
Radial only ($Z$) & no & 0.4944 & $[0.4666, 0.5273]$ \\
Tangential only ($X,Y$) & no & 0.4446 & $[0.4111, 0.4773]$ \\
\bottomrule\end{tabular}
\caption{Triaxial reconstruction, 10 adult participants, 84 trials each. The upper block is the complete $2\times3$ grid crossing partition type with filtering. The lower block reconstructs channel subsets and is not comparable to the full-array rows. Confidence intervals use participant bootstrap with 20{,}000 resamples and seed 0.}
\label{tab:s-tri-cohort}\end{table*}

\begin{table*}[htbp]\centering\small
\setlength{\tabcolsep}{4pt}
\begin{tabular}{lrrrrrrrr}
\toprule
Participant & Mono. raw & Mono. Soft & Rand. raw & Rand. Soft & Axis raw & Axis Soft & Radial & Tangential \\
\midrule
P01 & 0.3167 & 0.2813 & 0.3512 & 0.3206 & 0.4495 & 0.4094 & 0.4669 & 0.3837 \\
P02 & 0.3365 & 0.2991 & 0.3646 & 0.3334 & 0.4662 & 0.4290 & 0.4600 & 0.3869 \\
P03 & 0.3076 & 0.2734 & 0.3379 & 0.3077 & 0.4514 & 0.4136 & 0.4659 & 0.3736 \\
P04 & 0.3035 & 0.2698 & 0.3250 & 0.2932 & 0.4801 & 0.4348 & 0.4939 & 0.3746 \\
P05 & 0.4342 & 0.4038 & 0.4465 & 0.4275 & 0.5780 & 0.5530 & 0.5817 & 0.4800 \\
P06 & 0.4585 & 0.4369 & 0.4711 & 0.4523 & 0.5957 & 0.5744 & 0.5908 & 0.4744 \\
P07 & 0.4096 & 0.3948 & 0.4257 & 0.4052 & 0.5268 & 0.5113 & 0.4658 & 0.4911 \\
P08 & 0.3758 & 0.3560 & 0.4061 & 0.3890 & 0.5080 & 0.4908 & 0.4353 & 0.4907 \\
P09 & 0.3867 & 0.3667 & 0.4250 & 0.4062 & 0.5300 & 0.5102 & 0.4802 & 0.4865 \\
P10 & 0.3970 & 0.3796 & 0.4420 & 0.4183 & 0.5443 & 0.5228 & 0.5037 & 0.5048 \\
\midrule
Mean & 0.3726 & 0.3461 & 0.3995 & 0.3753 & 0.5130 & 0.4849 & 0.4944 & 0.4446 \\
\bottomrule\end{tabular}
\caption{Zero-shot reconstruction PCC per participant across the triaxial grid (Zenodo 11126593, the 10 adult participants). Every contrast in the main text is computed from these columns within a single filtering regime.}
\label{tab:s-tri-per-subject}\end{table*}

\begin{table*}[htbp]\centering\small
\begin{tabular}{lrrrr}
\toprule
Contrast (matched filtering) & Mean $\Delta$ & 95\% CI & Consistent & Exact $p$ \\
\midrule
\multicolumn{5}{l}{\textit{Grouping: axis partition against a size-matched random partition}} \\
Soft-HFC & $+0.1096$ & $[+0.1009, +0.1193]$ & 10/10 & 1.95e-03 \\
Raw & $+0.1135$ & $[+0.1040, +0.1253]$ & 10/10 & 1.95e-03 \\
\midrule
\multicolumn{5}{l}{\textit{Capacity: random partition against the monolithic tokenizer}} \\
Soft-HFC & $+0.0292$ & $[+0.0228, +0.0349]$ & 10/10 & 1.95e-03 \\
Raw & $+0.0269$ & $[+0.0205, +0.0333]$ & 10/10 & 1.95e-03 \\
\midrule
\multicolumn{5}{l}{\textit{Total: axis partition against the monolithic tokenizer}} \\
Soft-HFC & $+0.1388$ & $[+0.1313, +0.1467]$ & 10/10 & 1.95e-03 \\
Raw & $+0.1404$ & $[+0.1320, +0.1502]$ & 10/10 & 1.95e-03 \\
\midrule
\multicolumn{5}{l}{\textit{Filtering: Soft-HFC applied within each partition type}} \\
Monolithic & $-0.0265$ & $[-0.0314, -0.0215]$ & 10/10 & 1.95e-03 \\
Random 3-way & $-0.0242$ & $[-0.0278, -0.0206]$ & 10/10 & 1.95e-03 \\
Axis partition & $-0.0281$ & $[-0.0346, -0.0219]$ & 10/10 & 1.95e-03 \\
\bottomrule\end{tabular}
\caption{All paired triaxial contrasts, 10 participants. Every row compares two conditions with identical filtering, or one partition type under both filtering settings. With $n=10$ the smallest attainable exact $p$ is $2/2^{10}\approx0.002$, so values at that floor indicate perfect sign consistency rather than a resolved tail.}
\label{tab:s-tri-contrasts}\end{table*}

\begin{table*}[htbp]\centering\small
\setlength{\tabcolsep}{4pt}
\begin{tabular}{llrrrrr}
\toprule
Configuration & HFC & Probe & Ctrl A & Classical & Classical 95\% CI & Cl. $>$ Pr. \\
\midrule
Monolithic & no & 48.24 & 46.71 & 53.54 & $[49.93, 56.75]$ & 8/10 \\
Monolithic & Soft & 49.07 & 46.67 & 54.50 & $[51.03, 57.65]$ & 8/10 \\
Random 3-way & no & 50.83 & 48.69 & 53.54 & $[49.93, 56.75]$ & 7/10 \\
Random 3-way & Soft & 48.51 & 48.61 & 54.50 & $[51.03, 57.65]$ & 7/10 \\
Axis partition & no & 47.53 & 51.61 & 53.54 & $[49.93, 56.75]$ & 7/10 \\
Axis partition & Soft & 50.19 & 50.12 & 54.50 & $[51.03, 57.65]$ & 8/10 \\
Radial only ($Z$) & no & 47.19 & 47.79 & 53.33 & $[49.94, 57.01]$ & 8/10 \\
Tangential only ($X,Y$) & no & 49.35 & 46.78 & 52.76 & $[48.54, 56.75]$ & 8/10 \\
\bottomrule\end{tabular}
\caption{Braille decoding balanced accuracy (\%), five-fold within participant, 10 adult participants. The classical arm computes six log band powers spanning 1--96\,Hz plus four temporal window means per channel, so it does not depend on how channels are grouped for the tokenizer and takes the same value across partitions sharing a channel set and filtering. Its interval excludes chance only in the Soft-HFC full-array conditions. The final column counts participants in which the classical decoder exceeds the pretrained probe.}
\label{tab:s-decoding}\end{table*}



\clearpage
\bibliography{references}
\end{document}
